\documentclass[11pt,a4paper]{article}
\hbadness=10000   % Silences ALL "Underfull \hbox" warnings
\hfuzz=100pt      % Silences "Overfull \hbox" warnings up to 100 points
% ----------------------------------------------------------------------
%  Working knowledge base / executive record for the MSc project
%  "Visual Augmentation of Audio Semantics for Accessibility"
%  Maintained as a living document. Compiles on Overleaf or with pdflatex.
% ----------------------------------------------------------------------

\usepackage[utf8]{inputenc}
\usepackage[T1]{fontenc}
\usepackage[margin=1in]{geometry}
\usepackage{booktabs}
\usepackage{longtable}
\usepackage{enumitem}
\usepackage{xcolor}
\usepackage{titlesec}
\usepackage[T1]{fontenc}
\usepackage{lmodern}  % Scalable vector version of Computer Modern
\usepackage{microtype}
\usepackage[hidelinks]{hyperref}
\hypersetup{colorlinks=true, linkcolor=blue!50!black, urlcolor=blue!60!black, citecolor=blue!50!black}

\titleformat{\section}{\large\bfseries}{\thesection}{0.6em}{}
\setlist{nosep, leftmargin=1.4em}

% Inline action-item marker for things Adam must verify/decide.
% All such items are collected in Section "Action Items -- Tasks for Adam".
\newcommand{\TODO}[1]{\textcolor{red}{\textbf{[TODO--Adam: #1]}}}

\title{\textbf{Visual Augmentation of Audio Semantics for Accessibility}\\[4pt]
\large Working Knowledge Base: Literature Review, Design Notes \& Project Plan}
\author{Adam Gavriely \\ \small MSc Final Project}
\date{Last updated: \today}

\begin{document}
\maketitle

\begin{abstract}
\noindent This document is the living record of the MSc project. It summarises the refined
project definition, a grounded literature review across the eight areas named in the proposal,
concrete model and dataset recommendations with hardware implications, the evaluation design,
and a two-month work plan. It also logs the key decisions taken so far so that the full history
of the project is recoverable from a single place. Sections are updated as the work progresses;
a dated changelog is kept in Section~\ref{sec:changelog}.
\end{abstract}

\tableofcontents

% ======================================================================
\section{Executive Summary}
% ======================================================================
The project builds a system that makes the \emph{non-speech} information in a video's soundtrack
accessible to deaf and hard-of-hearing (DHH) viewers. Subtitles and sign-language interpretation
convey \emph{speech}, but they largely discard ambient and environmental sound --- flowing water,
approaching footsteps, a distant siren, rain, birds, a reacting crowd, machinery --- which carries
real meaning and atmosphere. The system detects this missing audio semantics and renders it as
\textbf{complementary visual augmentations displayed alongside the original video}, supplementing
(not replacing) the video and its captions.

The distinctive idea is \textbf{cross-modal gating}: because the \emph{video} is also analysed,
the system augments only what the audio conveys but the picture does \emph{not} already show.
A siren whose source is already on screen needs no augmentation; a siren from off-screen is
exactly what is worth illustrating. This ``reason about the gap between what is heard and what is
seen'' framing is the core contribution and the axis on which the method is compared to baselines
that ignore the visual stream.

No models are trained. The contribution is (i) the \emph{task formulation}, (ii) an inference
pipeline that orchestrates existing open-source foundation models, (iii) a curated benchmark, and
(iv) an automatic evaluation protocol. Work is scoped to roughly two months on a modest local GPU
(GTX~1060 6\,GB) with a university GPU requested for the diffusion and large-VLM stages.

\medskip
\noindent\textbf{Status at time of writing.} Project definition refined and locked; literature
review drafted (this document); model/dataset shortlist and hardware plan produced; two-month
schedule produced. Implementation of the new pipeline not yet started --- a prior code skeleton
(FFmpeg audio extraction, Whisper ASR wiring, modular stage architecture) is reusable.

% ======================================================================
\section{Project Definition and Scope}
% ======================================================================
\subsection{Problem statement}
Given a short video clip (10--20\,s) with synchronised audio and video, identify semantic
information in the audio --- with emphasis on ambient and environmental sound --- that is
\emph{insufficiently conveyed by the visual stream}, and generate complementary visual
augmentations that communicate it to a DHH viewer.

\subsection{Inputs and outputs}
\textbf{Input:} a short clip with synchronised video + audio. Candidate sources: movie/TV scenes,
documentaries, real-world recordings. The audio may contain speech, ambient sound, environmental
sound, discrete acoustic events, and background music.
\textbf{Output:} a set of complementary visual augmentations shown alongside the original video.
Static, storyboard-style images first; short generated video segments are a stretch goal.

\subsection{The seven-stage pipeline}
The proposal defines a modular inference pipeline in which each stage can be evaluated or replaced
independently:
\begin{enumerate}
  \item \textbf{Audio extraction} --- FFmpeg.
  \item \textbf{Video understanding} --- what is already visible (Qwen2.5-VL / LLaVA / InternVL).
  \item \textbf{Speech recognition} --- transcribe speech (Whisper Large V3 / Qwen2-Audio).
  \item \textbf{Audio event detection} --- non-speech sounds (BEATs / PANNs / AudioSet models).
  \item \textbf{Cross-modal semantic analysis} --- fuse scene context, speech and audio events;
        decide which cues matter, are \emph{not} already visible, and how to depict them
        (Qwen3 / Llama~3.1).
  \item \textbf{Visual augmentation generation} --- FLUX.1 / Stable Diffusion XL.
  \item \textbf{Evaluation}.
\end{enumerate}

\subsection{Scope boundaries and the ``no training'' constraint}
Off-the-shelf open-source foundation models are glued together; nothing is trained or fine-tuned.
The intellectual contribution is task formulation, cross-modal orchestration (especially stage~5),
benchmark construction, and evaluation methodology.

\subsection{Note on the change of direction}
An earlier version of this project illustrated the \emph{speech} of news clips for Deaf viewers.
The current proposal is a deliberate pivot: the star is now \emph{non-speech ambient/environmental
sound}, the \emph{video} stream is analysed (previously ignored), audio event detection and
cross-modal gating are new core components, output is overlaid alongside the original video, and a
formal 300-clip benchmark replaces ad-hoc news clips. Reusable from the earlier build: FFmpeg audio
extraction, Whisper ASR wiring, and the modular pluggable-stage architecture.

% ======================================================================
\section{Research Questions (from the proposal)}
% ======================================================================
\begin{itemize}
  \item Which types of audio information contribute most to scene understanding when presented visually?
  \item What level of semantic granularity should be used to represent ambient/environmental sounds?
  \item Can automatically generated visual augmentations improve accessibility beyond subtitles?
  \item How should semantic audio augmentation systems be evaluated?
  \item How does the proposed pipeline compare with representative audio-to-visual approaches?
\end{itemize}

% ======================================================================
\section{Literature Review}\label{sec:litreview}
% ======================================================================
This review covers the eight areas named in the proposal. It is deliberately organised around the
\emph{role each area plays in the pipeline} so that model choices in Section~\ref{sec:models}
follow directly from it.

% ----------------------------------------------------------------------
\subsection{Audio Event Detection and Tagging (Stage 4)}
% ----------------------------------------------------------------------
This is the engine that turns the soundtrack into a set of labelled non-speech sound events ---
the raw material for augmentation. The field is anchored by \textbf{AudioSet}~\cite{audioset}, a
2\,M-clip ontology of 527 sound classes that underpins nearly all modern systems. Two families
dominate:
\begin{itemize}
  \item \textbf{CNN-based tagging --- PANNs}~\cite{panns}. Large-scale pretrained audio neural
    networks (notably the CNN14 backbone) established strong general-purpose audio features and a
    then-state-of-the-art mAP of 0.439 on AudioSet tagging. PANNs remain a robust, lightweight,
    easy-to-run baseline and are ideal for a first implementation on a 6\,GB GPU.
  \item \textbf{Transformer-based tagging --- AST and BEATs}. The Audio Spectrogram Transformer
    (AST)~\cite{ast} applied a pure-transformer patch model to spectrograms. \textbf{BEATs}~\cite{beats}
    (audio pre-training with acoustic tokenizers) is the current practical state of the art and is
    used as the baseline encoder in the DCASE 2023/2024 Sound Event Detection challenges. For frame-
    level \emph{detection} (with time boundaries, not just clip-level tags), CRNN+BEATs and
    frame-level transformer variants (e.g.\ ATST-Frame) lead recent DCASE results.
\end{itemize}
\textbf{Relevance / choice.} We need both \emph{what} sound and \emph{when} (time boundaries) so
augmentations can be time-aligned. Recommendation: start with PANNs (CNN14) for fast, dependency-light
clip tagging; move to a BEATs-based tagger for quality; use a CRNN/frame-level model only if precise
onset/offset timing proves necessary. AudioSet's 527-class ontology also gives us a ready-made
vocabulary for the ``which sounds matter'' analysis in Stage~5.

% ----------------------------------------------------------------------
\subsection{Audio Captioning (Stage 4/5 alternative representation)}
% ----------------------------------------------------------------------
Where tagging outputs labels, \emph{automated audio captioning} (AAC) outputs a free-text sentence
describing the soundscape (e.g.\ ``a train approaches as birds chirp in the background''). This is a
richer, more compositional representation and a natural bridge to an image-generation prompt. The
canonical datasets are \textbf{AudioCaps}~\cite{audiocaps} (AudioSet-derived, $\sim$46\,k clips) and
\textbf{Clotho}~\cite{clotho}. Modern AAC systems pair a strong audio encoder with a language
decoder: \textbf{EnCLAP}~\cite{enclap} (neural-codec + CLAP + LLM) and \textbf{SLAM-AAC}~\cite{slamaac}
(current SOTA on Clotho/AudioCaps, using paraphrase augmentation and CLAP-based re-ranking) are
representative. \textbf{CLAP}~\cite{clap} (contrastive language--audio pretraining) is the workhorse
embedding underneath much of this and is directly useful to us as an audio$\leftrightarrow$text
relevance scorer in evaluation.
\textbf{Relevance / choice.} Captioning is attractive as the ``description'' that feeds image
generation, and it is a natural \textbf{baseline} (``audio captioning only, no visual output'').
It also provides the \emph{audio captioning} baseline the proposal requires.

% ----------------------------------------------------------------------
\subsection{Audio Understanding / Large Audio-Language Models (Stages 3--5)}
% ----------------------------------------------------------------------
Large Audio-Language Models (LALMs) fold ASR, audio-event reasoning, and captioning into a single
instructable model: an audio encoder + a modality connector + an LLM. Key systems: \textbf{Qwen2-Audio}
\cite{qwenaudio} (strong general audio understanding; can be prompted for both transcription and
sound description), \textbf{SALMONN}~\cite{salmonn} (speech + audio-event + music encoders on a
frozen LLM), and \textbf{Audio Flamingo~3}~\cite{audioflamingo3} (fully open, long-context,
multi-audio dialogue). Evaluation of these models is itself an active area; \textbf{MMAU}~\cite{mmau}
is a multi-task audio understanding/reasoning benchmark worth citing for how the community measures
these capabilities.
\textbf{Relevance / choice.} A LALM is an appealing way to \emph{collapse Stages 3--4} into one
component and to reason about audio directly in language. Trade-off: a single 7B+ LALM is heavier
and less individually inspectable than PANNs + Whisper as separate, swappable modules. Plan: keep
the modular ASR+AED path as the dependable default; treat a LALM (Qwen2-Audio) as an alternative
Stage-3/4 backend and a possible ablation (``LALM vs.\ specialised modules'').

% ----------------------------------------------------------------------
\subsection{Vision-Language and Multimodal LLMs (Stages 2 and 5)}
% ----------------------------------------------------------------------
Stage~2 (what is already visible) and Stage~5 (cross-modal fusion) both lean on vision-language
models. \textbf{Qwen2.5-VL}~\cite{qwenvl} is the leading open option: native dynamic-resolution
vision, explicit video understanding with temporal grounding (it can localise events to video
segments), and 3B/7B/72B sizes --- the 7B fits our hardware budget and the 3B is a fallback for the
local GPU. \textbf{InternVL3}~\cite{internvl} and the \textbf{LLaVA}~\cite{llava} family are the
principal alternatives. These same models double as the \textbf{VLM describer} in evaluation
(Section~\ref{sec:eval}).
\textbf{Relevance / choice.} Qwen2.5-VL-7B for video understanding and as the evaluation describer;
Qwen2.5-VL-3B or a frame-sampled call on the local GPU as fallback. Because the same model family
spans understanding \emph{and} evaluation, we should use \emph{different} model instances/prompts
for generation vs.\ judging to avoid a model grading its own output.

% ----------------------------------------------------------------------
\subsection{Audio-to-Image Generation (Stage 6 --- and a key baseline)}
% ----------------------------------------------------------------------
Two philosophies exist. (1) \textbf{Direct audio-to-image}: \textbf{Sound2Scene}~\cite{sound2scene}
learns to map audio into the latent space of a pretrained image generator and synthesises a scene
directly from sound, trained only on unlabelled video; \textbf{SonicDiffusion}~\cite{sonicdiffusion}
conditions a pretrained diffusion model on audio. These are exactly the proposal's ``direct
audio-to-image'' \textbf{baseline}. (2) \textbf{Text-conditioned generation} driven by a
language description of the sound: \textbf{Stable Diffusion XL}~\cite{sdxl} and
\textbf{FLUX.1}~\cite{flux} are the open state of the art. Our proposed method sits in camp (2) but
with the crucial extra step that the \emph{prompt is chosen by cross-modal reasoning} (Stage~5),
not produced blindly from audio.
\textbf{Relevance / choice.} Use SDXL / FLUX.1 for augmentation generation (text-conditioned on a
Stage-5 prompt); use Sound2Scene as the direct-audio-to-image baseline. This contrast --- reasoned,
gap-aware prompting vs.\ blind audio$\to$image --- is precisely what the evaluation must isolate.

% ----------------------------------------------------------------------
\subsection{Audio-to-Video Generation (Stage 6 stretch goal)}
% ----------------------------------------------------------------------
Extending static augmentations to short motion clips is the proposal's stretch. Directly
audio-conditioned video generation is still maturing; the pragmatic route is \emph{image-to-video}
applied to a generated keyframe (e.g.\ Stable Video Diffusion-style motion), keeping the reasoning
in Stages 2--5 and treating motion as a rendering choice. This keeps the stretch goal cheap and
optional and avoids depending on a fast-moving, heavyweight research area.

% ----------------------------------------------------------------------
\subsection{Accessibility Technologies for DHH Users (motivation \& evaluation grounding)}
% ----------------------------------------------------------------------
This literature justifies the whole project and shapes what ``useful'' means. Two threads matter:
\begin{itemize}
  \item \textbf{Limits of captioning for non-speech sound.} Standard captioning represents only a
    narrow slice of the soundtrack; substantial atmospheric and emotional information is lost, and
    non-speech sounds are inconsistently or poorly captioned~\cite{captionlimits}. Work on
    visualising prosody/emotion in captions~\cite{prosodycaptions} shows the community actively
    augmenting text with visual cues --- directly adjacent to our idea, but for speech rather than
    ambient sound.
  \item \textbf{Sound-awareness systems.} Jain et al.'s \textbf{SoundWatch}~\cite{soundwatch}
    (smartwatch sound classification/alerts) and \textbf{HomeSound}~\cite{homesound} (in-home sound
    visualisation, deployed with DHH households) establish, from field studies, \emph{which} sounds
    DHH users want to be aware of and how visual sound feedback is received. \textbf{ProtoSound}
    \cite{protosound} personalises sound recognition for DHH users. These ground our choice of
    ``which sounds matter'' and provide real user-centred motivation and evaluation framing.
\end{itemize}
\textbf{Relevance.} This is the section that separates the project from a generic multimodal demo:
it names the accessibility gap (non-speech sound), cites evidence that it is real and unmet, and
gives a principled basis for the ``which audio cues are worth showing'' decision in Stage~5.

% ----------------------------------------------------------------------
\subsection{Benchmark Datasets}\label{sec:datasets}
% ----------------------------------------------------------------------
The proposal names three sources for the $\sim$300-clip benchmark:
\begin{itemize}
  \item \textbf{VGGSound}~\cite{vggsound} --- $\sim$200\,k 10-second YouTube clips, 309 classes, each
    with a \emph{visible} sound source. Excellent for clean single-event ambient/environmental
    scenarios; its ``sound source is on screen'' property is useful both as positive material and,
    inverted, to construct ``sound source \emph{off} screen'' cases that stress cross-modal gating.
  \item \textbf{UnAV-100}~\cite{unav100} --- $\sim$10.8\,k untrimmed videos, 30\,k+ audio-visual
    events, 100 categories, avg 2.8 (often overlapping) events per video. Ideal for \emph{mixed
    audio scenes} and for cases with multiple concurrent sounds. CC-BY 4.0.
  \item \textbf{MovieNet}~\cite{movienet} --- 1{,}100 movies with rich annotations (42\,k scene
    boundaries, character/place/action/cinematic-style tags). Best source for cinematic
    \emph{acoustic-event} scenes (sirens, glass, explosions, applause) with narrative context.
\end{itemize}
\textbf{Selection strategy.} Sample across the three proposal scenarios --- ambient environmental
sound, localised acoustic events, mixed scenes --- prioritising clips where the audio carries
information \emph{only partially} present in the visuals (the property that makes the task
non-trivial). VGGSound for clean ambient/event clips, UnAV-100 for mixed/overlapping events,
MovieNet for narrative acoustic events. Record licence/source per clip.

\smallskip
\noindent\TODO{verify dataset access and terms before committing to the benchmark plan --- MovieNet
may require a signed agreement/application; VGGSound is distributed as YouTube links (some videos are
offline, so a download + availability check is needed); UnAV-100 is CC-BY 4.0 (fine).}
\TODO{confirm the 300-clip target is feasible within the timeline, or set a fallback (e.g.\ 100--150
clips).}

% ======================================================================
\section{The Accessibility Gap: Ambient Sound vs.\ Subtitles vs.\ Sign Language}\label{sec:gap}
% ======================================================================
This section is the motivational core of the project: \emph{what} non-speech audio information DHH
viewers miss, and \emph{how} the two dominant accessibility channels --- subtitles and sign language
--- handle (or fail to handle) it. It is the evidence base that justifies adding a visual channel.

\subsection{What ambient/non-speech sound do DHH users miss and want?}
Field studies of sound-awareness systems give a concrete, user-grounded answer. Jain et al.'s
\textbf{HomeSound}~\cite{homesound} and \textbf{SoundWatch}~\cite{soundwatch} were built around a set
of $\sim$19 everyday sounds ($\sim$31\,h of data) that DHH users said they wanted to know about:
door knocks/doorbells, a microwave or appliance beeping, a baby crying, a dog barking, a cat
meowing, running water, alarms, sirens, a running vehicle, birds chirping. Two observations matter
for this project:
\begin{itemize}
  \item \textbf{These are mostly \emph{alerting / functional / safety} sounds} --- things that
    prompt an action or signal danger in the user's real environment. That is a different need from
    \emph{media consumption}, where the missing sounds are often \emph{atmospheric and narrative}
    (rain setting a mood, an off-screen crowd, ominous music, distant thunder). Existing
    sound-awareness tech targets the former; this project targets the latter, comparatively
    under-served niche.
  \item \textbf{The taxonomy still transfers.} The categories --- environmental/ambient, discrete
    acoustic events, human/animal sounds --- line up with the proposal's three benchmark scenarios
    and with the AudioSet ontology used in Stage~4. \TODO{when defining the benchmark's ``which
    sounds matter'' labels, reuse this user-validated taxonomy rather than inventing one.}
\end{itemize}

\subsection{How subtitles / SDH handle non-speech sound}
The relevant caption type is \textbf{SDH} (Subtitles for the Deaf and Hard-of-hearing), the most
complete form, which --- unlike plain subtitles --- is expected to represent non-dialogue audio:
sound effects, music, speaker identity, manner of speaking, and audience reaction. Conventions are
codified in style guides such as the \textbf{DCMP Captioning Key}~\cite{dcmp}: sound effects are set
in square brackets, e.g.\ \texttt{[applause]}, \texttt{[glass shatters]}; music is marked with note
symbols. \emph{Its guidance is strikingly close to this project's core idea:} ``the source of the
sound effect should be included \dots \textbf{unless the source is clearly seen on screen}'' and
``if it can be inferred through the visuals, it is not necessary to caption it.'' In other words,
\textbf{professional captioners already apply cross-modal gating by hand} --- describe the sound only
when it is not already visible. Stage~5 of this project \emph{automates a principle the captioning
industry endorses but performs manually}. This is a strong justification to foreground.

\medskip\noindent\textbf{Where SDH falls short (the opportunity):}
\begin{itemize}
  \item \textbf{Low bandwidth / lost granularity.} A rich soundscape is compressed to a terse text
    label. A University of Sheffield study found that in \emph{Jaws}, captions told viewers there was
    ``famous music'' but \emph{not} that the music signalled the shark approaching --- so DHH viewers
    lost the suspense the score was built to create~\cite{sheffield}.
  \item \textbf{Emotional/atmospheric information is lost.} An EEG study found that attending only to
    standard compliant captions loses substantial emotional information, and that non-verbal-sound
    captions tend to trigger extra \emph{attentional} effort rather than an emotional
    response~\cite{captionlimits}.
  \item \textbf{Divided attention.} Reading captions competes with watching the scene; DHH viewers
    split visual attention between two sources, causing fatigue and missed on-screen content (visual
    gags, who is speaking)~\cite{livecaptions}. \TODO{design risk to validate: our augmentations are
    \emph{also} visual, so they must not \emph{worsen} divided attention --- placement/timing/salience
    of augmentations is a real design question and a candidate evaluation dimension.}
  \item \textbf{No spatial or continuous information.} SDH cannot easily convey \emph{where} a sound
    is (off-screen left? approaching?) or how it evolves --- exactly the information a visual
    augmentation can carry.
  \item \textbf{Inconsistency.} Non-speech captioning quality varies by provider/regulation; recent
    HCI work is actively trying to enrich it (visualizing speech style~\cite{prosodycaptions,
    speechstyles}; adapting non-speech captions generatively --- CapTune~\cite{captune}), which shows
    the community sees the same gap this project addresses.
\end{itemize}

\subsection{How sign language handles non-speech sound}
Sign language can carry \emph{more} than subtitles: skilled interpreters convey music and sound
through classifiers, mimed sound sources (air-guitar, an instrument), sustained holds for long
notes, rhythm, and \textbf{non-manual markers} (facial expression for pitch/intensity), aiming to
transmit the \emph{feeling} of the sound, not just its name --- as seen in theatrical and concert
interpreting. Limits for this project's setting:
\begin{itemize}
  \item \textbf{Availability.} Interpretation is live, human, and costly; the vast majority of
    recorded/broadcast media has no interpreter, and where present it centres on \emph{speech} (and
    sometimes music), rarely systematic \emph{ambient/environmental} sound.
  \item \textbf{Automation is unsolved.} Automatic sign-language \emph{generation} remains an open
    research problem (noted in the proposal's related work), so it is not a deployable channel for
    ambient sound today.
  \item \textbf{Takeaway.} Sign language shows the \emph{ceiling} --- rich, expressive, affective
    conveyance of sound --- that terse SDH text does not reach. A generated \emph{visual}
    augmentation is a pragmatic, automatable middle path: higher-bandwidth and more atmospheric than
    a bracketed label, without requiring a human interpreter or solved sign-language synthesis.
\end{itemize}

\subsection{Synthesis --- why this project}
DHH viewers miss real semantic and atmospheric information in non-speech sound. Subtitles/SDH capture
it only as low-bandwidth text that loses granularity, emotion, spatial and temporal nuance, and adds
reading/attention load; and they already \emph{manually} decide to describe a sound only when it is
not visible --- the very gating this project automates. Sign language conveys sound richly but is
unavailable at scale and not automatable for ambient sound. This positions automatically generated,
gap-aware \emph{visual augmentation} as a complementary channel that is (a) higher-bandwidth than
SDH text, (b) automatable unlike interpretation, and (c) grounded in a principle captioners already
endorse. \TODO{frame the report's introduction around this three-way comparison --- it is the
clearest articulation of the contribution.}

% ======================================================================
\section{Model and Tooling Recommendations}\label{sec:models}
% ======================================================================
Shortlist per stage, with a ``local-first'' option that fits the GTX~1060 (6\,GB) and a ``quality''
option for the university GPU.

\begin{center}
\begin{longtable}{@{}p{2.7cm}p{4.3cm}p{4.6cm}p{2.4cm}@{}}
\toprule
\textbf{Stage} & \textbf{Local-first (6\,GB) / dev} & \textbf{Quality (uni GPU)} & \textbf{Licence} \\
\midrule
\endhead
1. Audio extract & FFmpeg & FFmpeg & permissive \\
2. Video understanding & Qwen2.5-VL-3B (frame sampling) & Qwen2.5-VL-7B / InternVL3 & Apache-2.0* \\
3. Speech recognition & faster-whisper (small/medium) & Whisper Large V3 & MIT/Apache \\
4. Audio events & PANNs (CNN14) & BEATs tagger & MIT / research \\
5. Cross-modal analysis & Llama~3.1-8B / Qwen3-8B (quant.) & Qwen3-14B+ or hosted API & varies \\
6. Generation & SD 1.5 / SDXL-Turbo (test) & SDXL / FLUX.1-dev & CreativeML/Apache* \\
7. Evaluation & CLAP + CLIP scores & VLM describe + LLM judge & --- \\
\bottomrule
\end{longtable}
\end{center}
\footnotesize *Verify each model's exact licence at integration time; FLUX.1 has \texttt{dev}
(non-commercial) vs.\ \texttt{schnell} (Apache-2.0) variants, and some audio-event checkpoints are
research-only. \normalsize
\smallskip

\noindent\TODO{confirm the university GPU VRAM once granted --- the split between ``local-first'' and
``quality'' columns above depends on it.} \TODO{confirm an acceptable FLUX.1 licence variant
(\texttt{dev} non-commercial vs.\ \texttt{schnell} Apache-2.0) for an academic deliverable; check
with the supervisor whether a non-commercial model is acceptable.}

\medskip
\noindent\textbf{Hardware reality check.} The GTX~1060 (6\,GB) can run FFmpeg, faster-whisper
(small/medium), PANNs, and SD~1.5 / SDXL-Turbo for development, plus quantised small LLMs slowly.
It \emph{cannot} comfortably run SDXL/FLUX at quality, a 7B VLM, or a 14B+ LLM. \textbf{Action:
request the university GPU early} (lead time) for the quality runs; do all wiring and the benchmark
build locally in the meantime. A hosted API for the Stage-5 LLM is a viable fallback if the GPU is
delayed.

% ======================================================================
\section{Evaluation Methodology}\label{sec:eval}
% ======================================================================
\subsection{Automatic protocol (primary)}
Following the proposal:
\begin{enumerate}
  \item Analyse the original video + audio.
  \item Generate the complementary visual augmentation(s).
  \item A VLM \emph{describes} the augmented scene (video + augmentation).
  \item Derive a \emph{semantic reference} from the original multimodal input via an LLM --- i.e.\
    what a hearing viewer would take from the soundtrack that is not visible.
  \item An independent \textbf{LLM judge} scores the semantic consistency between the augmentation's
    described content and the reference.
\end{enumerate}
\textbf{Design cautions.} (a) Use \emph{different} models/prompts for generation and judging to
avoid self-preference bias. (b) Complement the LLM judge with cheaper, less gameable signals:
\textbf{CLAP} audio$\leftrightarrow$image and \textbf{CLIP}~\cite{clip} text$\leftrightarrow$image
relevance. (c) Report the judge's agreement with a small human-labelled subset to validate it.

\subsection{Metrics}
Semantic-consistency score (LLM judge, rubric-based); augmentation relevance (CLIP/CLAP);
\emph{gating accuracy} --- did the system correctly \emph{withhold} augmentation when the sound is
already visible, and add it when not? This last metric is the direct measure of the core
contribution and should be built on a hand-labelled subset.

\subsection{Baselines (required by the proposal)}
(1) Direct audio-to-image (Sound2Scene) --- no visual analysis; (2) Audio captioning only --- text,
no visual; (3) Proposed cross-modal method. The comparison isolates the value of reasoning about the
audio--visual gap.

\subsection{Human evaluation (optional)}
If time allows: a small qualitative study comparing ``subtitles alone'' vs.\ ``subtitles +
augmentation'' on scene-understanding questions. Given the two-month scope and ethics/recruitment
overhead, treat this as optional and keep the automatic protocol as the primary evidence.
\smallskip

\noindent\TODO{decide whether to attempt the human study at all; if yes, check the department's
ethics/IRB requirements and lead time early --- recruitment/ethics can be the schedule bottleneck.}

% ======================================================================
\section{Project Plan (approx.\ 8 weeks)}\label{sec:plan}
% ======================================================================
Assumes $\sim$2 months, local GTX~1060 for development, university GPU requested in Week~1. Weeks are
indicative and overlap; the benchmark build (data-gathering) runs in parallel with implementation.

\begin{center}
\begin{longtable}{@{}p{1.5cm}p{6.4cm}p{5.2cm}@{}}
\toprule
\textbf{Week} & \textbf{Focus} & \textbf{Deliverable / milestone} \\
\midrule
\endhead
0 (now) & Refine scope; literature review; plan; \textbf{request uni GPU}; set up fresh repo structure.
  & This document; GPU request submitted; repo scaffold. \\
1 & Stage 1--3 skeleton (FFmpeg, faster-whisper, and a first VLM call for video understanding);
    reuse old wiring. & End-to-end stub runs on one clip. \\
2 & Stage 4 audio event detection (PANNs first, then BEATs); produce labelled+timed events. &
    Sound events extracted on sample clips. \\
3 & Stage 5 cross-modal analysis: the gap-aware ``what to augment'' component (LLM prompt design). &
    First real augmentation prompts; qualitative check. \\
4 & Stage 6 generation (SDXL/FLUX on uni GPU) + compositing augmentations alongside video. &
    First full augmented clips end-to-end. \\
5 & \textbf{Benchmark build}: curate $\sim$100 (scaling to 300) clips from VGGSound/UnAV-100/MovieNet
    across the three scenarios; record licences. & Benchmark v1 + selection criteria documented. \\
6 & Evaluation harness: VLM-describe + LLM-judge + CLIP/CLAP; implement gating-accuracy metric;
    wire the two baselines. & Automatic evaluation runs; baseline numbers. \\
7 & Full experiments across the benchmark; ablations (modular vs.\ LALM; with/without cross-modal
    gating); error analysis. & Results tables + qualitative examples. \\
8 & Write-up: technical report, polish prototype, finalise deliverables. &
    Report draft; reproducible code; curated benchmark. \\
\bottomrule
\end{longtable}
\end{center}

\noindent\textbf{Critical path \& risks.} (1) University GPU lead time --- request in Week~0; hosted
API fallback for Stage~5. (2) Dataset access/download volume (VGGSound is large; download only
sampled clips). (3) Scope creep on the human study and audio-to-video --- both explicitly optional.
(4) Evaluation validity --- validate the LLM judge against a small human-labelled subset early.

% ======================================================================
\section{Action Items --- Tasks for Adam}\label{sec:tasks}
% ======================================================================
Open questions and external actions that only Adam can resolve. These are the inline
\TODO{}\ markers, collected. Tick them off as they are settled and note the outcome in
Section~\ref{sec:changelog}.

\begin{enumerate}[label=\textbf{A\arabic*.}]
  \item \textbf{Compute --- submit the university GPU request now.} Lead time is the schedule's top
    risk. Confirm the granted \textbf{VRAM} on approval; it decides which models run locally vs.\ on
    the cluster (Section~\ref{sec:models}).
  \item \textbf{Dataset access \& terms.} MovieNet may need a signed agreement/application; VGGSound
    is YouTube links (check availability, plan downloads); UnAV-100 is CC-BY 4.0 (fine). Verify
    before committing to the benchmark plan (Section~\ref{sec:datasets}).
  \item \textbf{Benchmark size.} Confirm $\sim$300 clips is feasible in the timeline, or agree a
    fallback (100--150) with the supervisor.
  \item \textbf{Model licences.} Confirm an acceptable \textbf{FLUX.1} variant (\texttt{dev}
    non-commercial vs.\ \texttt{schnell} Apache-2.0); check specific \textbf{audio-event checkpoints}
    (some BEATs/DCASE checkpoints are research-only).
  \item \textbf{Human study decision.} Decide whether to run the optional qualitative study; if yes,
    check department ethics/IRB requirements and lead time early.
  \item \textbf{Scope of ``background music''.} The proposal lists music as a possible audio
    component but the focus is ambient/environmental sound. Confirm with the supervisor whether music
    is in or out of scope for augmentation.
  \item \textbf{Stages 3--4 design choice.} After a small pilot, decide the default: modular
    (Whisper + PANNs/BEATs) vs.\ a single LALM (Qwen2-Audio). Keep the other as an ablation.
  \item \textbf{Reference details.} Before the final report, verify the flagged bibliography entries
    (exact author lists/venues) --- see \TODO{} in Section~\ref{sec:litreview}/References.
  \item \textbf{Security carry-over.} Confirm the previously leaked GitHub PAT was revoked/regenerated
    (flagged in an earlier session).
\end{enumerate}

% ======================================================================
\section{Decisions and Notes Log}\label{sec:changelog}
% ======================================================================
\begin{itemize}
  \item \textbf{2026-08-06 --- Project pivot confirmed.} Definition changed from ``illustrate the
    speech of news clips'' to ``augment video with visuals of missing non-speech audio semantics,
    gated by video understanding.'' New definition is authoritative.
  \item \textbf{2026-08-06 --- Codebase decision.} Fresh start, reusing pieces (FFmpeg audio,
    Whisper wiring, modular architecture) from the earlier skeleton.
  \item \textbf{2026-08-06 --- First deliverables.} Literature review + project plan (this document).
  \item \textbf{2026-08-06 --- Accessibility deep-dive added} (Section~\ref{sec:gap}): what ambient
    sound DHH users miss, and how SDH subtitles and sign language handle it. Key finding: SDH already
    performs \emph{manual} cross-modal gating (don't caption a visible sound source), which Stage~5
    automates. Full source trail in the Research Log (Section~\ref{sec:researchlog}).
  \item \textbf{2026-08-06 --- Constraints unchanged.} $\sim$2 months; local GTX~1060 6\,GB;
    university GPU to be requested.
  \item \textbf{Open items.} Confirm university GPU VRAM once granted; finalise exact benchmark
    composition after a first pass over the datasets; decide modular-vs-LALM default for Stages 3--4
    after a small pilot.
\end{itemize}

% ======================================================================
% Note moved outside the environment to prevent list parsing errors
% ======================================================================
\noindent\TODO{verify the exact author lists/venues/years of the entries marked ``[verify]'' below before the final report --- these were assembled from search results and not all are fully confirmed.}\par\smallskip

\begin{thebibliography}{99}

\bibitem{audioset} J.\ F.\ Gemmeke et al., ``AudioSet: An ontology and human-labeled dataset for audio events,'' \emph{ICASSP}, 2017.
\bibitem{panns} Q.\ Kong et al., ``PANNs: Large-Scale Pretrained Audio Neural Networks for Audio Pattern Recognition,'' \emph{IEEE/ACM TASLP}, 2020. arXiv:1912.10211.
\bibitem{ast} Y.\ Gong, Y.-A.\ Chung, J.\ Glass, ``AST: Audio Spectrogram Transformer,'' \emph{Interspeech}, 2021. arXiv:2104.01778.
\bibitem{beats} S.\ Chen et al., ``BEATs: Audio Pre-Training with Acoustic Tokenizers,'' \emph{ICML}, 2023. arXiv:2212.09058.
\bibitem{audiocaps} C.\ D.\ Kim et al., ``AudioCaps: Generating Captions for Audios in The Wild,'' \emph{NAACL}, 2019.
\bibitem{clotho} K.\ Drossos et al., ``Clotho: An Audio Captioning Dataset,'' \emph{ICASSP}, 2020.
\bibitem{clap} B.\ Elizalde et al., ``CLAP: Learning Audio Concepts from Natural Language Supervision,'' \emph{ICASSP}, 2023.
\bibitem{enclap} J.\ Kim et al., ``EnCLAP: Combining Neural Audio Codec and Audio-Text Joint Embedding for Automated Audio Captioning,'' \emph{ICASSP}, 2024. [verify authors]
\bibitem{slamaac} W.\ Chen et al., ``SLAM-AAC: Enhancing Audio Captioning with Paraphrasing Augmentation and CLAP-Refine through LLMs,'' 2024. arXiv:2410.09503.
\bibitem{qwenaudio} Y.\ Chu et al., ``Qwen-Audio / Qwen2-Audio: Advancing Universal Audio Understanding,'' 2023--2024. arXiv:2311.07919.
\bibitem{salmonn} C.\ Tang et al., ``SALMONN: Towards Generic Hearing Abilities for Large Language Models,'' \emph{ICLR}, 2024.
\bibitem{audioflamingo3} ``Audio Flamingo 3: Advancing Audio Intelligence,'' 2025. arXiv:2507.08128. [verify authors]
\bibitem{mmau} S.\ Sakshi et al., ``MMAU: A Massive Multi-Task Audio Understanding and Reasoning Benchmark,'' 2024. arXiv:2410.19168.
\bibitem{qwenvl} Qwen Team, ``Qwen2.5-VL Technical Report,'' 2025. arXiv:2502.13923.
\bibitem{internvl} ``InternVL3: Native Multimodal Pretraining,'' 2025. [verify authors/venue]
\bibitem{llava} H.\ Liu et al., ``Visual Instruction Tuning (LLaVA),'' \emph{NeurIPS}, 2023.
\bibitem{sound2scene} K.\ Sung-Bin et al., ``Sound to Visual Scene Generation by Audio-to-Visual Latent Alignment (Sound2Scene),'' \emph{CVPR}, 2023. arXiv:2306.11504.
\bibitem{sonicdiffusion} ``SonicDiffusion: Audio-Driven Image Generation and Editing with Pretrained Diffusion Models,'' 2024. [verify authors]
\bibitem{sdxl} D.\ Podell et al., ``SDXL: Improving Latent Diffusion Models for High-Resolution Image Synthesis,'' 2023. arXiv:2307.01952.
\bibitem{flux} Black Forest Labs, ``FLUX.1,'' 2024. \url{https://blackforestlabs.ai/}.
\bibitem{captionlimits} ``Limitations of Standard Accessible Captioning of Sounds and Music for DHH People: An EEG Study,'' \emph{Frontiers in Integrative Neuroscience}, 2020. [verify authors]
\bibitem{prosodycaptions} ``Visualization of Speech Prosody and Emotion in Captions: Accessibility for Deaf and Hard-of-Hearing Users,'' \emph{CHI}, 2023.
\bibitem{soundwatch} D.\ Jain et al., ``SoundWatch: Smartwatch-based Deep Learning Approaches to Support Sound Awareness for DHH Users,'' \emph{ASSETS}, 2020.
\bibitem{homesound} D.\ Jain et al., ``HomeSound: An Iterative Field Deployment of an In-Home Sound Awareness System for DHH Users,'' \emph{CHI}, 2020.
\bibitem{protosound} D.\ Jain et al., ``ProtoSound: A Personalized and Scalable Sound Recognition System for DHH Users,'' \emph{CHI}, 2022. arXiv:2202.11134. \url{https://arxiv.org/abs/2202.11134}
\bibitem{dcmp} Described and Captioned Media Program (DCMP), ``Captioning Key --- Guidelines and Preferred Techniques'' (Sound Effects, Music, Non-Speech Information sections). \url{https://dcmp.org/learn/captioningkey}
\bibitem{sheffield} University of Sheffield / SubTxt Creative, study on captions and suspense in film and TV, 2024. \emph{[verify: locate the primary publication; currently via project coverage]} \url{https://liamodell.com/2024/05/09/captions-cinema-suspense-university-of-sheffield-paper-research-project-survey-subtxt-creative/}
\bibitem{livecaptions} ``How Users Experience Closed Captions on Live Television: Quality Metrics Remain a Challenge,'' 2024. arXiv:2404.10153. \url{https://arxiv.org/abs/2404.10153} [verify authors]
\bibitem{captune} ``CapTune: Adapting Non-Speech Captions With Anchored Generative Models,'' 2025. arXiv:2508.19971. \url{https://arxiv.org/abs/2508.19971} [verify authors]
\bibitem{speechstyles} ``Visualizing Speech Styles in Captions for Deaf and Hard-of-Hearing Viewers,'' \emph{Int.\ J.\ Human-Computer Studies}, 2024. \url{https://www.sciencedirect.com/science/article/abs/pii/S1071581924001691} [verify authors]
\bibitem{vggsound} H.\ Chen et al., ``VGGSound: A Large-scale Audio-Visual Dataset,'' \emph{ICASSP}, 2020. arXiv:2004.14368.
\bibitem{unav100} T.\ Geng et al., ``Dense-Localizing Audio-Visual Events in Untrimmed Videos (UnAV-100),'' \emph{CVPR}, 2023. arXiv:2303.12930.
\bibitem{movienet} Q.\ Huang et al., ``MovieNet: A Holistic Dataset for Movie Understanding,'' \emph{ECCV}, 2020. arXiv:2007.10937.
\bibitem{clip} A.\ Radford et al., ``Learning Transferable Visual Models From Natural Language Supervision (CLIP),'' \emph{ICML}, 2021.
\bibitem{whisper} A.\ Radford et al., ``Robust Speech Recognition via Large-Scale Weak Supervision (Whisper),'' 2022. arXiv:2212.04356.
\end{thebibliography}

% ======================================================================
\appendix
\section{Research Log --- Sources Consulted and Conclusions}\label{sec:researchlog}
% ======================================================================
A record of the background research behind this document: the topics investigated, the key primary
links, and the conclusion drawn from each. Web searches were run in August 2026; where a claim rests
on secondary coverage rather than a confirmed primary paper it is tagged for verification.

\footnotesize
\subsection*{Audio event detection (Stage 4)}
Key links: PANNs \url{https://arxiv.org/abs/1912.10211}; AST \url{https://arxiv.org/abs/2104.01778};
BEATs \url{https://arxiv.org/abs/2212.09058}; AudioSet tagging code
\url{https://github.com/qiuqiangkong/audioset_tagging_cnn}.
\textbf{Conclusion:} PANNs/CNN14 is the light, dependable first implementation; BEATs is the quality
upgrade and the DCASE baseline encoder; use a frame-level model only if precise onset/offset timing
is required. AudioSet's 527-class ontology is our sound vocabulary.

\subsection*{Audio captioning (baseline + prompt source)}
Key links: SLAM-AAC \url{https://arxiv.org/abs/2410.09503}; AudioSetCaps
\url{https://arxiv.org/abs/2411.18953}.
\textbf{Conclusion:} captioning gives a compositional soundscape description usable as a generation
prompt and doubles as the required ``audio captioning only'' baseline; CLAP is reusable as an
audio$\leftrightarrow$text relevance scorer in evaluation.

\subsection*{Large audio-language models (Stages 3--4 alternative)}
Key links: Qwen-Audio/Qwen2-Audio \url{https://arxiv.org/abs/2311.07919}; Audio Flamingo 3
\url{https://arxiv.org/abs/2507.08128}; MMAU benchmark \url{https://arxiv.org/abs/2410.19168}.
\textbf{Conclusion:} a LALM can collapse ASR+AED into one instructable component; keep modular
Whisper+PANNs as the inspectable default and treat the LALM as an ablation.

\subsection*{Vision-language models (Stages 2 and 7)}
Key links: Qwen2.5-VL \url{https://arxiv.org/abs/2502.13923}; Qwen2.5-VL-7B weights
\url{https://huggingface.co/Qwen/Qwen2.5-VL-7B-Instruct}.
\textbf{Conclusion:} Qwen2.5-VL-7B for video understanding and as the evaluation describer (3B as
local fallback); it has explicit temporal grounding. Use different instances/prompts for generation
vs.\ judging.

\subsection*{Audio-to-image / diffusion (Stage 6 + direct baseline)}
Key links: Sound2Scene \url{https://sound2scene.github.io/} (code
\url{https://github.com/kaist-ami/Sound2Scene}); SonicDiffusion
\url{https://cyberiada.github.io/SonicDiffusion/}; SDXL \url{https://arxiv.org/abs/2307.01952};
FLUX.1 \url{https://blackforestlabs.ai/}.
\textbf{Conclusion:} generate with text-conditioned SDXL/FLUX driven by the Stage-5 (gap-aware)
prompt; use Sound2Scene as the ``direct audio-to-image, no visual reasoning'' baseline --- the
contrast is exactly what the evaluation isolates.

\subsection*{Benchmark datasets}
Key links: VGGSound \url{https://arxiv.org/abs/2004.14368} / \url{https://github.com/hche11/VGGSound};
UnAV-100 \url{https://unav100.github.io/} / \url{https://arxiv.org/abs/2303.12930}; MovieNet
\url{https://arxiv.org/abs/2007.10937} / \url{https://movienet.github.io/}.
\textbf{Conclusion:} VGGSound for clean ambient/event clips (visible source; invert for
``off-screen'' cases), UnAV-100 for mixed/overlapping events (CC-BY 4.0), MovieNet for narrative
acoustic events. Access/terms need checking (Task~A2).

\subsection*{Accessibility --- what DHH users miss and how subtitles/sign language cope}
Key links: SoundWatch \url{https://dl.acm.org/doi/10.1145/3373625.3416991}; HomeSound
\url{https://soundability.eecs.umich.edu/img/portfolio/Jain_HomeSound_CHI2020.pdf}; DCMP Captioning
Key \url{https://dcmp.org/learn/captioningkey}; EEG caption-limitation study
\url{https://www.ncbi.nlm.nih.gov/pmc/articles/PMC7040021/}; live-caption UX
\url{https://arxiv.org/abs/2404.10153}; CapTune \url{https://arxiv.org/abs/2508.19971}; Sheffield
suspense/captions coverage
\url{https://liamodell.com/2024/05/09/captions-cinema-suspense-university-of-sheffield-paper-research-project-survey-subtxt-creative/}.
\textbf{Conclusions:} (1) DHH users' desired sounds are mostly \emph{alerting/functional}
(SoundWatch/HomeSound); media \emph{atmospheric} sound is the under-served niche this project targets.
(2) SDH is the only caption type that covers non-speech sound, but it is low-bandwidth text that loses
granularity (the \emph{Jaws} suspense example), emotion (EEG study), and spatial/temporal nuance, and
adds divided-attention load. (3) The DCMP rule ``don't caption a sound whose source is visible on
screen'' means professional captioners already perform \emph{manual cross-modal gating} --- Stage~5
automates an endorsed principle. (4) Sign language conveys sound richly but is unavailable at scale
and not automatable for ambient sound, so it marks the expressive ceiling rather than a deployable
channel. Together these motivate gap-aware visual augmentation as a complementary channel.
\normalsize

\end{document}